% =========================================================
\section{Giai Đoạn Phân Tích và Kỹ Nghệ Dữ Liệu}
% =========================================================

Sau khi yêu cầu được xác nhận, hệ thống phân tích cấu trúc cơ sở dữ liệu và chuẩn bị đầu vào cho mô hình hóa đồ thị. Giai đoạn này gồm ba tác tử hoạt động tuần tự, mỗi tác tử sinh một sản phẩm (artifact) mà tác tử kế tiếp sử dụng.

\subsection{Tác Tử Phân Tích ($\mathcal{A}_{eda}$)}

$\mathcal{A}_{eda}$ khám phá và trích xuất siêu dữ liệu từ cơ sở dữ liệu quan hệ. Đây là tác tử đầu tiên tương tác trực tiếp với dữ liệu thô; kết quả phân tích của nó định hình các quyết định thiết kế ở toàn bộ giai đoạn sau.

\paragraph{Quy trình phân tích.}
$\mathcal{A}_{eda}$ thực thi tuần tự ba pha qua tám công cụ chuyên dụng. \textit{Pha chuẩn bị} xác thực kết nối cơ sở dữ liệu, xuất toàn bộ bảng sang định dạng phẳng, và trích xuất siêu dữ liệu lược đồ (khóa chính, khóa ngoại, cột thời gian). \textit{Pha phân tích} thực hiện bốn phép phân tích: (i)~thống kê mô tả cho từng bảng và cột, (ii)~phát hiện vấn đề chất lượng dữ liệu theo ba mức nghiêm trọng, (iii)~xác định các cột thời gian và đề xuất ngưỡng chia tập val/test dựa trên phân vị thời gian, và (iv)~phân loại quan hệ giữa các bảng theo vai trò --- bảng dữ kiện (fact) hay bảng chiều (dimension) --- tương ứng với phân loại trong Mục~\ref{sec:relational_data}. \textit{Pha tổng hợp} gộp kết quả thành báo cáo toàn cảnh phục vụ mô hình hóa.

Kết quả được lưu thành \textbf{bản đặc tả lược đồ} $\Phi$ trong $\mathcal{S}_{schema}$, xác định vai trò của từng bảng: bảng thực thể chính, bảng ngữ cảnh, và cột thời gian tham chiếu. Phân loại fact/dimension này hướng dẫn các bước tiếp theo: bảng thực thể (dimension) được dùng làm đối tượng dự đoán, bảng sự kiện (fact) cung cấp ngữ cảnh lịch sử trong truy vấn SQL.

\subsection{Tác Tử Xây Dựng Dữ Liệu ($\mathcal{A}_{builder}$)}

$\mathcal{A}_{builder}$ tự động hóa bước chuyển đổi từ cơ sở dữ liệu quan hệ $(\mathcal{T}, \mathcal{L})$ sang đồ thị dị thể $G=(\mathcal{V},\mathcal{E},\phi,\psi)$ đã định nghĩa trong Mục~\ref{sec:relation-graph}. Sản phẩm đầu ra tuân thủ giao diện \texttt{Dataset} của RelBench~\cite{robinson2024relbench}.

\paragraph{Nguyên tắc ánh xạ.}
Quá trình xây dựng đồ thị tuân theo ba quy tắc: (i)~mỗi hàng trong bảng $T_i$ trở thành một nút $v \in \mathcal{V}$ với kiểu $\tau_v = T_i$, định danh bởi khóa chính; (ii)~mỗi quan hệ khóa ngoại giữa $T_i$ và $T_j$ trở thành cạnh có hướng, kèm cạnh ngược tương ứng với $\mathcal{L}^{-1}$ (Mục~\ref{sec:schema-graph}) để đảm bảo truyền thông điệp hai chiều; (iii)~mỗi cột được gán bộ mã hóa phù hợp với kiểu dữ liệu theo hệ thống mã hóa đa phương thức trong Mục~\ref{sec:encoders} --- các cột khóa được loại bỏ khỏi tập đặc trưng.

\paragraph{Xác thực mã sinh ra.}
Mã do $\mathcal{A}_{builder}$ sinh ra phải vượt qua cổng kiểm soát chất lượng trước khi được ghi vào tệp: kiểm tra cú pháp, xác minh sự tồn tại của lớp giao diện bắt buộc, và phát hiện thư viện thiếu khai báo import. Mã không hợp lệ bị từ chối kèm phản hồi lỗi cụ thể, theo nguyên tắc \textit{fail-fast}.

\subsection{Tác Tử Xây Dựng Tác Vụ ($\mathcal{A}_{task}$)}

$\mathcal{A}_{task}$ tự động hóa quá trình xây dựng \textbf{bảng huấn luyện} $T_{train}$ với cấu trúc $(\mathcal{K}_v, t_v, y_v)$ đã định nghĩa trong Mục~\ref{sec:task_to_training_table}. Tác tử sinh truy vấn SQL qua DuckDB để tạo bảng $\mathcal{D}_{train} = \{(e_i, y_i, t_i)\}$, trong đó $e_i$ là thực thể dự đoán, $y_i$ là nhãn tính từ dữ liệu lịch sử, và $t_i$ là mốc thời gian tham chiếu.

\paragraph{Quy trình xây dựng.}
$\mathcal{A}_{task}$ tuân theo quy trình bốn bước cố định, quy định tường minh trong system prompt: (1)~xác định loại tác vụ từ độ đo đánh giá hoặc mô tả mục tiêu; (2)~xác thực ràng buộc thời gian (khoảng cách giữa dấu thời gian val và test phải không nhỏ hơn cửa sổ dự đoán $\Delta$); (3)~phân tích cấu trúc tác vụ để chọn mẫu SQL phù hợp; (4)~sinh, kiểm thử, và đăng ký mã. Sau khi sinh mã, hệ thống xác minh loại tác vụ nhúng trong mã trùng khớp với ý định người dùng --- nếu không khớp, tệp bị xóa và tác tử thử lại, ngăn lỗi lan truyền đến giai đoạn huấn luyện.

\paragraph{Hệ thống mẫu thiết kế SQL.}
System prompt nhúng năm mẫu thiết kế SQL, mỗi mẫu phù hợp với một cấu trúc dữ liệu và ngữ nghĩa tác vụ khác nhau. Việc lựa chọn mẫu dựa trên phân tích tự động cung cấp: nguồn gốc thực thể, thống kê khoảng cách thời gian giữa các sự kiện, và tín hiệu ngữ nghĩa. Bảng~\ref{tab:sql-patterns} tóm tắt năm mẫu.

\begin{table}[htbp]
\centering
\caption{Năm mẫu thiết kế SQL cho bảng huấn luyện và điều kiện lựa chọn.}
\label{tab:sql-patterns}
\begin{tabular}{lp{4cm}p{6.5cm}}
\toprule
\textbf{Mẫu} & \textbf{Ngữ nghĩa tác vụ} & \textbf{Điều kiện lựa chọn} \\
\midrule
\textbf{A} & Churn, vắng mặt hành vi & Bảng thực thể tồn tại; kiểm tra sự tồn tại sự kiện trong cửa sổ $\Delta$ \\
\textbf{B} & Sự kiện có khoảng cách lớn & Phân vị 90 khoảng cách sự kiện $> \Delta$; cửa sổ lookback xác định từ dữ liệu \\
\textbf{C} & Dự đoán toàn thực thể & Bảng thực thể tồn tại; giá trị 0 hợp lệ (ví dụ: 0 giao dịch) \\
\textbf{D} & Thực thể có ngày tạo & Cột ngày tạo phát hiện được; lọc thực thể theo thời điểm xuất hiện \\
\textbf{Link} & Dự đoán liên kết & Tác vụ link prediction; trả về danh sách thực thể đích \\
\bottomrule
\end{tabular}
\end{table}

Quyết định giữa mẫu A và B phụ thuộc vào phân vị 90 khoảng cách sự kiện: khi giá trị này vượt $\Delta$ (ví dụ: sự kiện theo mùa, giao dịch theo quý), kiểm tra sự tồn tại trong cửa sổ $\Delta$ sẽ bỏ sót thực thể trong khoảng trống, và mẫu B với cửa sổ lookback dài hơn là lựa chọn phù hợp.

\paragraph{Phòng chống rò rỉ dữ liệu.}
Mốc thời gian $t_i$ trong bảng huấn luyện kích hoạt chiến lược lấy mẫu lân cận nhất quán theo thời gian (Mục~\ref{sec:time-consistent}): chỉ các cạnh và nút xuất hiện trước $t_i$ được sử dụng khi huấn luyện, ngăn chặn rò rỉ dữ liệu từ tương lai. Ngưỡng chia tập val/test được xác định từ phân phối thời gian do $\mathcal{A}_{eda}$ phân tích.
